\usepackage[utf8]{inputenc}
\usepackage[spanish,es-tabla]{babel}
\usepackage{amsmath,amssymb,amsfonts}
\usepackage{graphicx}
\usepackage{geometry}
\usepackage{booktabs}
\usepackage{cite}
\usepackage{caption}
\usepackage{abstract}
\usepackage{titlesec}
\usepackage[colorlinks=true, linkcolor=blue, urlcolor=blue, citecolor=blue]{hyperref}

% --- MÁRGENES ---
\geometry{
	letterpaper,
	top=2.0cm,
	bottom=2.0cm,
	left=1.8cm,
	right=1.8cm,
	columnsep=0.6cm
}

% --- FORMATO DE SECCIONES ---
\titleformat{\section}
{\normalfont\scshape\centering}
{\thesection.}
{0.5em}
{}

\renewcommand{\thesection}{\Roman{section}}

\titleformat{\subsection}
{\normalfont\itshape}
{\thesubsection.}
{0.5em}
{}

\renewcommand{\thesubsection}{\Alph{subsection}}

% --- CONFIGURACIÓN DEL RESUMEN ---
\renewcommand{\abstractname}{}
\setlength{\absleftindent}{0pt}
\setlength{\absrightindent}{0pt}

% Desactiva el título principal del resumen
\renewcommand{\abstractname}{}
\renewcommand{\absnamepos}{empty}
% Eliminar el título "Resumen" del entorno abstract por completo
\makeatletter
\renewenvironment{abstract}{%
	\small
	\par\noindent
}{%
	\par
}
\makeatother